%%
%%

\subsection{Referenzaufgaben}

\subsubsection{Referenzaufgabe: Rechtsseitiger Hypothesentest}
Ein Hersteller von Kanülen behauptet, dass maximal 0.1\,\% seiner Verpackungen beschädigt seien. Sie wollen einen Test mit 1000 Stück durchführen.

Erstellen Sie einen Hypothesentest auf einem Signifikanzniveau von $\alpha \le 1\,\%$.

\TNTeop{
  Eine Tabelle mit $n=1000, p=0.001$ liefert:

  \begin{tabular}{cc}
    $k$ & $P(X <= k)$ \\\hline
    0   & 36.77\,\% \\\hline
    1   & 73.58\,\% \\\hline
    2   & 91.98\,\% \\\hline
    3   & 98.11\,\% \\\hline
    4   & 99.64\,\% \\\hline
    5   & 99.94\,\% \\\hline
    6   & ...\,\%\\\hline
  \end{tabular}
  
  Naiv könnte man auf die Idee kommen bei $k=3$ haben wir das linksseitige Intervall gefunden. Doch wenn wir nun das Intervall $[0;3]$ als Annahmebereich festlegen, wäre $[4;1000]$ unser Ablehnungsbereich.
  der Annhmebereich von $[0;3]$ hat aber eine Wahrscheinlichkeit von 98.11\,\%, was dem Ablehnungsbereich eine Wahrscheinlichkeit von $1-0.9811$, also von 1.9\% gibt. Damit ist unser $k=3$ noch zu klein. Da binomialcdf immer von links rechnet, müssen wir noch eins dazunehmen, damit wir im Ablehnungsbereich definitiv unter 1\,\% fallen.

  Unser $k$ ist 4, der Annahmebereich ist $[0;4]$ und der Ablehnungsbereich liegt bei $[5;1000]$. Wenn also mehr als vier Verpackungen defekt sind, haben wir Grund zur Annahme, dass etwas an den 0.1\,\% nicht stimmt.
  }

\begin{rezept}{Rechtsseitiger Ablehnungsbereich}{}
  Beim rechtsseitigen Ablehnungsbereich ist typischerweie unser $k$ eins höher zu wählen, da \texttt{binomial\textbf{c}df} immer von links her addiert.

  Im Zweifelsfalle machen Sie die Probe: Wie viel Prozent liegt denn nun im Ablehnungsbereich?
\end{rezept}
\newpage

